\documentclass[11pt]{article}
\usepackage[margin=1in]{geometry}
\usepackage{amsmath,amssymb,amsfonts,amsthm,mathtools,bm}
\usepackage[T1]{fontenc}
\usepackage{microtype}
\newtheorem{theorem}{Theorem}
\newtheorem{lemma}{Lemma}
\newtheorem{cor}{Corollary}
\newtheorem{defn}{Definition}
\newtheorem{principle}{Principle}
\newtheorem{statement}{Statement}
\theoremstyle{remark}
\newtheorem{statementr}{Statement}

\title{The Bohmian Story: Mechanics, Measurement, and Objections}
\author{James Taylor}
\date{Working draft derived from the 2003 thesis}
\begin{document}
\maketitle

\begin{quote}\small This is a thesis-derived expository draft. It preserves the original presentation as a baseline; its historical discussion and references will need updating before publication.\end{quote}

\section{A brief word about Bohmian Mechanics}


Briefly, Bohmian mechanics is a quantum theory which considers the
state of an $N$-particle system to be the wave function and the
locations of the $N$ particles. The dynamics is specified by Schr\"odinger's equation\ and
a first-order equation for the evolution of the positions. Simply put,
the particles are guided by the wave which is why wave behavior is
seen. The particle behavior occurs because there are
particles. Although not an obvious fact, it is true that the 
 entire quantum measurement formalism
can be derived from an analysis of Bohmian mechanics \cite{DGZ92, DGZ03}.
 One of the advantages of 
 Bohmian mechanics is that it is easy to understand the starting
point. There are no axioms about measurements. The whole theory is the
dynamical system mentioned above. Although Bohmian mechanics is
different from classical mechanics as its laws of motion are
different, it is not a philosophically radical departure from the
starting point of mechanics. That is to say, it is still based on a particle
ontology.

\clearpage

\section{Bohmian mechanics} \label{s:sim}


Bohmian mechanics is a theory about 
 particles with definite locations. The fundamental objects to specify
 are  the trajectories in physical space of these particles. 
The object which
 defines the trajectories is 
the wave function, familiar from quantum
 mechanics. More precisely,  the state of the system in Bohmian mechanics\ is given by the
 pair $(Q, \ensuremath{\psi})$; $Q = 
(\boldsymbol Q_1, \ldots, \boldsymbol Q_N) \in \ensuremath{\mathbb{R}^{\ensuremath{3} N}}$
is the configuration of the $N$ particles in our system  and \ensuremath{\psi}\ is
a  wave function\ 
on the configuration space \ensuremath{\mathbb{R}^{\ensuremath{3} N}}, as usual in a quantum theory.
The state of the system changes according to Bohm's equation and Schr\"odinger's equation: 
\begin{equation} \label{e:bohm}
\ensuremath{\frac{d{\boldsymbol{Q}_{k}}}{d{t}}} = \ensuremath{\frac{\hbar}{m_{k}}} \textrm{Im}
\frac {\ensuremath{({\ensuremath{\psi}}, {\nabla_{k} \ensuremath{\psi}})}}{ \ensuremath{({\ensuremath{\psi}}, {\ensuremath{\psi}})}}
(\boldsymbol{Q}_1,\ldots, \boldsymbol{Q}_N) \ensuremath{=:} v_k^{\ensuremath{\psi}} (Q) 
 \quad  k= 1,\ldots,N
\end{equation} 
\begin{equation}
\ensuremath{i \hbar \ensuremath{\frac{\partial {\ensuremath{\psi}}}{\partial {t}}} = - \sum_{k =
1}^{N} \ensuremath{\frac{\hbar^2}{2m_{k}}} \ensuremath{\Delta}_{k} \ensuremath{\psi} + V \ensuremath{\psi}   } = H\ensuremath{\psi},
\end{equation}
where $V$ is the potential function and \ensuremath{({\phi}, {\ensuremath{\psi}})} is the 
inner product on the value space, which we call the local inner
product, in distinction from the inner product on the Hilbert space of
wave functions. For complex-valued wave functions, the
potential is  a real-valued
function on configuration space and the inner product is 
$\bar{\phi}{\ensuremath{\psi}}$, where the bar denotes complex conjugation. 


\section{Spinor Bohmian mechanics}

To incorporate spin, the story barely changes.  
The value space for the wave function\  changes
from \ensuremath{\mathbb{C}}\ to 
the spinor space, \ensuremath{W}. By assumption, for one particle, \ensuremath{W}\ is a
complex, irreducible representation space of $SU(2)$, the covering
group of the rotation group $SO(3)$. If it is the spin $s$
representation  then $\ensuremath{W} = \ensuremath{\ensuremath{\mathbb{C}}^{2s +1}}$. The representation of the group
$SU(2)$ leads to a representation of the group's Lie algebra which is
the same Lie algebra as that of $SO(3)$. 
The operators   \ensuremath{{S_{x}}}, \ensuremath{{S_{y}}},
\ensuremath{{S_{z}}} are the three generators of the representation that
correspond to the 
infinitesimal rotations about the $x$-axis, $y$-axis, and $z$-axis,
respectively. They
form the spin vector  $\ensuremath{\boldsymbol{S}}_s = (\ensuremath{{S_{x}}}, \ensuremath{{S_{y}}},
\ensuremath{{S_{z}}})_s$. 
For $N$ particles, the value space becomes $ \ensuremath{\mathbb{C}}^{2s_1 +1}
\otimes \cdots \otimes \ensuremath{\mathbb{C}}^{2s_N +1}$,


 where the $i^{th}$ particle has spin $s_i$.  
The spin vector for each particle, $\ensuremath{\boldsymbol{S}}^{(i)}:= \ensuremath{\text{Id}} \otimes
\ldots \otimes  \ensuremath{\text{Id}}
\otimes S_{s_i} \otimes \ensuremath{\text{Id}} \otimes \ldots \otimes \ensuremath{\text{Id}} $, is the 1-particle
spin vector for spin $s_i$ acting only on the
 $i^{th}$ factor: 
\begin{equation*} 
\ensuremath{\boldsymbol{S}}^{(i)} (w_1 \otimes \cdots \otimes w_N)
\ensuremath{:=}
w_1 \otimes \cdots \otimes w_{i-1} \otimes \ensuremath{\boldsymbol{S}}_{s_i} (w_i) \otimes
w_{i+1} \otimes \cdots \otimes w_{N}.
\end{equation*}
The potential $V$ takes the form 
\begin{equation*}
U(q) + \sum_{i = 1}^{N} \ensuremath{\mu}_i \ensuremath{\boldsymbol{S}}^{(i)} \cdot \boldsymbol{B}(\ensuremath{\boldsymbol{q}}_i)
\end{equation*}
where $U$ is a real-valued function and $\boldsymbol{B}$ is the
magnetic field. Bohm's equation is unchanged, except we now interpret
\ensuremath{({\psi}, {\phi})} as a spinor inner product. The elements of the spin
spaces are spinors, a terminology we use even in the case of
the many particle spin space. 


\section[The quantum story]{The
emergence of the 
quantum formalism}

In this section, we shall give a rough explanation
for how the quantum formalism emerges
from Bohmian mechanics. The full explanation may be found in
\cite{DGZ92, DGZ03}.
The collapse of the wave function\ is the simplest part
of the story. The hardest part, and the one which we shall not say
very much about, is the emergence of the correct empirical statistics
in measurements. Somewhere in between in difficulty is the
understanding of why self-adjoint operators are related to
measurements. The discussion in this section is only for spinless
particles. 

We start with the notion of the conditional wave function. The wave function\ in Bohmian
mechanics is the wave function\ for the entire universe. What about wave functions for a
subsystem? A subsystem is some collection of particles. We shall use
$X$ for the system's actual configuration and use $x$ as a variable on
the configuration space of those particles. The particles not involved in
the subsystem will be called the environment and we shall use $Y$ for
its actual configuration and $y$ as the variable. The wave function\ of the
universe, $\Psi$, may then be written in terms of those two variables: 
$\Psi(x,y)$. The Bohmian velocity for the universal configuration is
$v^{\Psi} (X,Y)$. To obtain a wave function\ for the subsystem, we define the
conditional wave function\ $\ensuremath{\psi}(x) \ensuremath{:=} \Psi(x,Y)$.\footnote{For particles with
  spin, R. Tumulka has recently pointed out that density matrices
  replace the conditional wave function.} Immediately, one sees that
the Bohmian velocity of the conditional wave function\ at $X$ agrees with the
actual Bohmian velocity. The conditional wave function's evolution is governed
by the evolution of $\Psi$ as well as the evolution of the environment
$Y$. In certain situations, the conditional wave function\ will evolve according
to a Schr\"odinger\ equation defined on the subsystem; we then say that the
conditional wave function\ is an
effective wave function. In other situations, the
conditional wave function\ will not evolve according to Schr\"odinger's equation. 

In a
measurement situation, the system under consideration is assumed to
have an effective wave function.
 During a measurement, the
environment's evolution is influenced by the evolution of the
system. This generally will lead to the system's wave function\ no longer
evolving according to a Schr\"odinger\ equation. Instead, it evolves in a
completely different manner. This is the process of collapse. After
the measurement is over, the conditional wave function\ often returns to
being an effective wave function. This is how collapse occurs in Bohmian
mechanics. In terms of the actual Bohmian evolution of the universe,
nothing special happens when collapse occurs. Furthermore, collapse is
not a precise statement; although the conditional wave function\ exists and
changes as it changes, the length of collapse is presumably the
period of time in which conditional wave function\ is not
considered to be an effective wave function. This, of course, depends on the
level of approximation desired in the evolution of the conditional
wave function.  

In quantum mechanics, self-adjoint operators are related to
measurements. What do these very abstract mathematical objects have to
do with measurements? Mathematically, self-adjoint operators
bijectively correspond to projection-valued measures (PVM); this is
the spectral theorem of functional analysis. The measure is a measure
on the results of the experiments. They encapsualte the relevant
behavior of the environment. The projection implements the collapse of
the wave function. This is true for repeatable measurements. In more general
measurements, such as when the measurement cannot be repeated, 
the projections are replaced
with positive operators.   

There remains the question of the probabilities. That is to
say, we need to explain why $\ensuremath{\lvert {\ensuremath{\psi}} \rvert}^2$ gives the distribution of
the results of measurements. 
In short, 
the empirical agreement between Bohmian mechanics\ and standard quantum mechanics is
grounded in equivariance. 
  In Bohmian mechanics,  if the 
configuration is initially randomly distributed according to $|\ensuremath{\psi}_0|^2$,
then the evolution is such that the configuration will 
be distributed according to $|\ensuremath{\psi}_t|^2$ at time $t$. 
This is
equivariance. More precisely, the measure $\ensuremath{\rho}_0 = 
|\ensuremath{\psi}_0|^2$ can evolve either according to the evolution of \ensuremath{\psi}\ under
the Schr\"odinger\ evolution or as a density evolving according to the Bohmian
flow. Equivariance is the statement that those two evolutions agree.  
The continuity equation, a consequence of Schr\"odinger's equation, is 
\begin{equation}
\ensuremath{\frac{\partial {\ensuremath{\rho}}}{\partial {t}}} + \nabla \cdot (\ensuremath{\rho} v^{\ensuremath{\psi}}) = 0 
\end{equation}
and it guarantees equivariance.  

This answer 
is rather lacking in usefulness, even though it may seem satisfying. 
 Bohmian mechanics is a deterministic
dynamics. In particular, the initial
configuration of the universe is whatever it is. A satisfactory
statement, and a true statement,
 would be that for a typical initial configuration of the
universe, the universe evolves in such a way that the experimental
results of quantum mechanics emerges. Typicality is with respect to
$\ensuremath{\lvert {\Psi_0} \rvert}^2$. This is what is explained in \cite{DGZ92} and I
cannot find a short way of explaining it appropriately. 


\section{Particular experiments}


Any theory, no matter how beautiful its mathematics or sound its
explanations in abstract terms, must give a convincing explanation of
important experiments. We collect here some of the most prominent
experiments in quantum mechanics and explain the picture as presented
by Bohmian mechanics. We emphasize that this is technically unnecessary
as the experimental information is encoded in \ensuremath{|\ensuremath{\psi}|^2}. Thus the
quantum analysis of the system is sufficient even from a Bohmian
perspective. Indeed, the point of the
quantum formalism, as seen from the Bohmian perspective, is as a
thermodynamical-type formalism which allows one to predict the outcome of
experiments without having to do a detailed analysis of the, rather
complicated, Bohmian motions. Nevertheless, it is important to
appreciate what the measurements are or, more appropriately, are not 
measuring. 

All of the experiments below
are described in very rough terms; our intention is to only give a
rough story of the situation and not a detailed analysis beyond the
need for an informative and plausible argument. 

\subsection{A measurement of ``momentum''}

Let us consider a free particle. The experiment is the following. We
start with a wave function\ essentially supported around a position
$\boldsymbol{q}_0$; this 
can be created by
a nondestructive measurement of position. We then let
the system evolve and measure the position at some later time
$t$. A wave packet, with wabe vector $\boldsymbol k$, will evolve
so that its center is approximately $\boldsymbol{q}_0 +
\frac{\ensuremath{\hbar}}{m}\boldsymbol k 
t$, i.e.\ if $\ensuremath{\psi}_0 \approx A(\boldsymbol{q}) e^{{i}
  \boldsymbol k \cdot \boldsymbol{q}}  $ 
with $A(\boldsymbol{q})$ slowly varying for most of its support and
the support 
being relatively small centered around $\boldsymbol{q}_0$, 
then it will approximately evolve to $\ensuremath{\psi}_t \approx A_t(\boldsymbol{q})
e^{{i} \boldsymbol k \cdot \boldsymbol{q}} $ where 
$A_t$ has a somewhat larger support with  center 
$\boldsymbol{q}_t \ensuremath{:=} (\boldsymbol{q}_0
+\frac{\hbar}{m}\boldsymbol k t)$. If we do a position measurement at time
$t$, we
shall find the particle in the packet i.e.\ near $q_t$. Thus,
 it might seem natural to regard as the particle's momemtum 
 $m \frac{\boldsymbol{q}_t -
   \boldsymbol{q}_0}{t} = \hbar \boldsymbol{k} $. For
a general wave function\ we then consider a superposition of such  initial wave
packets. Given enough time, they will separate into different wave
packets whose centers are quite distant. The particle will be
distributed according to the norm squared of the coefficients in the
original superposition. That is to say, we shall see the distribution
of the Fourier transform when we do repeated measurements. Notice that
all we needed to do in order to measure momentum in this sense 
is to be able to
measure the approximate position of the particle.


Is this a measurement of the initial momentum of the actual Bohmian
particle? Absolutely not. There is no reason to believe that the
particle which ends up in the wave packet for momentum $\hbar \boldsymbol k$
actually has that as an initial momentum. Additionally, we do not
have conservation of the quantity $m \ensuremath{\boldsymbol{v}}$.  


\subsection{Scattering experiments} 

The above momentum measurements lead to a nice explanation of  
 scattering
experiments. Shoot at a system of interest a particle in
 a wave packet, one with a
prescribed momentum $\hbar \boldsymbol{k}_0$. It will interact
with the system and then separate from the system and again proceed
 freely. By setting up detectors, one can find the position of the
 particle and use the momentum measurement idea to deduce
 $\hbar \boldsymbol{k}_t$. By repeating this, one can find the distribution
 of the Fourier transform of the resultant wave function\ after interacting with
 the system.

What does this tell us about the scatterer and what does it do to the
scatterer? To answer this, we should write down the Hamiltonian of the
combined system. We have the Hamiltonian of the target, $H_o$, and of
the projectile, $H_p$. There also needs to be an interaction term
$H_{int}$ which should be localized so that it is significant 
only when the projectile is
near the target. The projectile's
Hamiltonian shall be assumed to be free. We wish to determine what
$V_o$ is. Let us assume that $H_o$ has eigenstates and that $\psi_0$
is a product of an eigenstate of $H_o$ with a wave packet of momentum
$\hbar \ensuremath{\boldsymbol{k}}$ directed towards the interaction region. That is to say, assume
 $\psi_0 = \lvert E_M \rangle \otimes \lvert \phi_{\ensuremath{\boldsymbol{k}}} \rangle$
 where $H_o \lvert E_M \rangle = E_M
\lvert E_M \rangle$. Let $H$ be the Hamiltonian of the system, i.e.\ 
$ H\ensuremath{:=} H_o +H_p +H_{int}$.


We then have that $H
\psi_0 \approx E_M \psi_0
 -\frac{\ensuremath{\hbar}^2 \ensuremath{\boldsymbol{k}} \cdot \ensuremath{\boldsymbol{k}}}{2m} \psi_0$ where the approximation is coming from
 assuming that the wave packet is essentially a plane wave on its
 support and that the initial wave function\ is such that $H_{int}$ is
 irrelevant.  
The evolution is such that the
 packet travels into the target, interacts with it, and then leaves the
 interaction region. 
 Let $U_t$ be the unitary evolution operator; thus  $H U_t \psi_0 = U_t
 H \psi_0$. When the wave function\ is supported away from the interaction
 region, then $H \approx H_o + H_p$. Thus, if $\ensuremath{\psi}_t$ is such that
 this is a good approximation, then $(H_o +H_p) U_t \psi_0 = U_t
 (H_o +H_p) \psi_0$ and we see that 
 $\psi_t$ is  an approximate
 eigenstate of the system with the same eigenvalue. Thus, $H \psi_t \approx
 (E_M -\frac{\ensuremath{\hbar}^2 \ensuremath{\boldsymbol{k}} \cdot \ensuremath{\boldsymbol{k}}}{2m} )\psi_t$. After the
 interaction with the object, the state can be represented
 approximately as $\psi_t \approx \sum_{n, \ensuremath{\boldsymbol{l}}} C_{n, \ensuremath{\boldsymbol{l}}} 
\lvert E_n \rangle \otimes
 \lvert \phi_{\ensuremath{\boldsymbol{l}}, t} \rangle$ where $\phi_{\ensuremath{\boldsymbol{l}}}$ is a wave
 packet  has
 momentum $\ensuremath{\hbar} \ensuremath{\boldsymbol{l}}$, evolves under the free evolution, and
 originates at the target's location.
 Applying $H$ to this state, we find $H \psi_t
 \approx \sum_{n, \ensuremath{\boldsymbol{l}}} (E_n - \frac{\ensuremath{\boldsymbol{l}} \cdot \ensuremath{\boldsymbol{l}}}{2}) C_{n, \ensuremath{\boldsymbol{l}}} \lvert E_n \rangle \otimes
 \lvert \phi_{\ensuremath{\boldsymbol{l}}, t} \rangle$. Since the terms in the sum are
 approximately orthogonal (approximately 
disjoint supports in configuration space
 because of the wave packets), we find that $E_M -\frac{\ensuremath{\hbar}^2\ensuremath{\boldsymbol{k}} \cdot \ensuremath{\boldsymbol{k}}}{2m} = E_n - \frac{\ensuremath{\hbar}^2\ensuremath{\boldsymbol{l}} \cdot \ensuremath{\boldsymbol{l}}}{2m}$ for every term in the sum such that $C_{n, \ensuremath{\boldsymbol{l}}} \neq 0$. That is to say, the eigenvalue differences in the initial
 and final state of the target are equal to the
 differences in the energy of the projectile, 
as provided by a momentum measurement.
  By using the
 superposition principle, one sees that this measurement procedure
 will only yield the eigenvalue differences; the differences are the only
 physically relevant information since the ground state energy
 corresponds to an irrelevant choice of a constant in the potential. 

This explains the
 discrete values obtained when we measure energy for systems with bound
 states. In particular, thinking of photons as particles, if we have
 various atoms and have photons interacting with them, then 
 the energy differences in the photons lie in a prescribed set of values
 which depends on the particular type of atom that they are interacting
 with.


As a final note about scattering experiments, observe that in Bohmian
mechanics we may also analyze time-of-arrival experiments. Such
experiments do not quite fit the quantum formalism, but physicists, of
course, do
have methods for obtaining the correct predictions.
An analysis of such
experiments may be found in \cite{DGTeZ00}.

\subsection{The double-slit experiment}


The double slit experiment is one of the most famous experiments in
all of quantum mechanics. It is often taken to be the example of all
the quantum mysteries. Bohmian mechanics accounts for the phenomena in a
trivial manner.

The setup is that of particles being directed towards a system with
two slits and a position detector, e.g.\ a screen,  
placed a distance away from the
slits. The system is arranged so that the wave function\ is approximately a plane
wave when it hits the two slits. After it passes through the slits,
there will be interference from the two emerging waves. The interference
pattern can be seen on the screen, built up out of many dots. 
If the particles are sent one at a
time, this pattern is built up over time. This is often considered to be
rather mysterious. Bohmian mechanics explains
that the dots on the screen occur because there are particles hitting
the screen; the pattern develops as a result of  the
particles being guided by a wave. 

If one of the slits is closed, then the interference pattern
disappears. This follows since there is now only one wave and thus
there is no interference, ignoring diffraction. The particle, being
guided by the wave, will also not reflect any interference. 

There is a
final modification of this experiment and this is supposed to be the
most troubling aspect of the experiment. 
Place a detector in such a way that the slit
which the particle goes through is identified, but arrange it so that,
classically at least, one would expect no disruption in the motion of
the particle. It happens that there is a disruption in the
motion. The reason is that the wave function\ is not  a wave in physical space,
but rather it is a wave in configuration space; detection changes the
relevant configuration space. 
In this experiment, as in all experiments, we are
using the effective wave function\ of the system. This is the actual wave function\ of the
entire universe with the irrelevant degrees of freedom conditioned
on. Furthermore, we are assuming that this conditioned wave function\ is
evolving according to Schr\"odinger's equation for the system. When we detect
the position of the particle,  more degrees of freedom become
relevant; in particular, the positions of the particles in the
detector. This separates the two parts of the projectile's wave function. That
is, the conditional wave function\ ceases to be an effective wave function. We can no
longer use the Schr\"odinger\ evolution of the system to evolve the original wave function. But,
once we take into account that the two parts of the wave function\ have
separated and that only one portion is relevant, the one with the
actual particle as reflected in the environment, 
then we can return to evolving that portion according to Schr\"odinger's
equation for the system. The conditional wave function\ automatically takes into
account the separation and collapse; all that one needs to check is
that the conditional wave function\ is effective. 


\subsection{Schr\"odinger's cat}

Picking up from the end of the double slit experiment discussion, we
can ask about Schr\"odinger's cat. The experiment is  to hook up a
quantum system to the fate of a cat. That is, do some measurement of
some quantum mechanical quantity, such as a spin measurement, and
decide the fate of a cat based on the outcome of that experiment. The
dilemma comes from the axioms of quantum mechanics. When is the cat's
fate decided? Does the collapse of the system occur only when someone
with a PhD comes along \cite{Bell87}? Or can a cat collapse itself because
it is alive or because it is a big system? 


Bohmian mechanics resolves it by stating that the cat has a definite
state because there are actual particles with definite positions. But
predicated upon this is that the state is recognizable. This really
means that the wave function's support decomposes into two, essentially disjoint,
supports. This is where size matters. Indeed, imagine a
cat which is alive versus one which is dead. Just by breathing, having
the blood flowing, wagging the tail, etc.,  many particles are in a different
position than they would be if the cat was dead. Let us say that the
average separation of each particle in the live cat versus the dead cat is
about 1 cm. Making the assumption that there about $10^{23}$ particles
in the cat, we find that the separation distance in configuration space
is roughly $\sqrt{10^{23}}$ cm; this is a rather large distance. 
One would therefore expect that the configurations corresponding to the
live or dead cat remain separated.
 That is, we have an effective collapse. 

The cat is illustrative of a different side of collapse. For a
microscopic system, it is the entanglement of the environment which
causes the collapse. But if one enlarges the system to include the
relevant environmental degrees of freedom, then the collapse does not
happen. Instead, one has a very large separation between the different
possibilties. It is the location of the particle which then determines
the relevant packet. But it is only an approximation to discard the
other parts of the wave function. For the conditional wave function, there is no
approximation except for when to consider it an effective wave function. For
macroscopic systems, the system itself evolves in such a way that
portions of the wave function\
will have negligible influence on the configuration's evolution.


\subsection{Position measurements}

We now describe a measurement of the position operator which does not
accurately measure the position of the particle. By a measurement of
the operator, we mean an experiment whose statistical information is
encapsulated in the associated PVM. In the case of position, given a
wave function\ \ensuremath{\psi}, a measurement of position is one in which the result $R$ is
distributed according to $\ensuremath{\lvert {\ensuremath{\psi}} \rvert}^2$ as the measurement is repeated
on identically prepared systems. 


The measurement described in this section will satisfy that
 criterion, but whether physicists 
 would accept the measurement or not as a measurement of
position is not clear to this author. But, from a Bohmian perspective,
 it is definitely not a measurement of position. Bohmian mechanics
provides a way of actually determining whether the experiment did
measure what it was supposed to have measured. 

Before discussing the measurement, we need to mention some important
facts about Bohmian mechanics. It is a first-order theory. This
implies that trajectories do not cross. We have already mentioned that
it is equivariant. In one dimension, this means that if $q_0$ is the
initial starting point, then $q_t$ is, more or less, uniquely defined by 
\[
\int_{-\infty}^{q_0} \ensuremath{\lvert {\ensuremath{\psi}_0} \rvert}^2 (x) dx = 
\int_{-\infty}^{q_t} \ensuremath{\lvert {\ensuremath{\psi}_t} \rvert}^2 (x) dx  .
\] 
As a particular application,  if
$\ensuremath{\lvert {\ensuremath{\psi}_t} \rvert}^2$ is periodic in time, then the Bohmian velocities are
periodic in time with the same period. We emphasize that this is
special to one dimension.


Our example is in one dimension with the harmonic oscillator potential $V(x) =
x^2$, ignoring constants. The evolution of any wave function\ 
is fully periodic in time with 
period $4 \pi$ and periodic up to a sign in $2\pi$ units of time. 
In $\pi$ units of time, the wave function\ $\psi$ will be $i \psi(-x)$. Let $x_m$ be
the median of the probability distribution $\ensuremath{\lvert {\psi} \rvert}^2$.  If a particle
starts to the left of the median, then it remains to the left of the median. 
Thus, if it begins to the left of $x_m$, 
then at time $t = \pi$, it is to the left of
the median at that time which is $-x_m$. 

If we look at the probability
distribution, we will see that it has been reflected across the
origin. The position operator at $t = \pi$, $\hat{Q}_{\pi}$, is
therefore the same as minus the position operator at $t=0$, 
$- \hat{Q}_{0}$. 
This means that we can measure the initial position operator by
measuring the position operator at time $t=\pi$ and then taking
minus that result as the result for measuring the initial position 
operator. If we assume that the measurement of the operator at time
$t=\pi$ accurately measures the position of the particle at that time,
then the above measurement for $\hat{Q}_0$ disagrees with the actual
initial position except if the particle initially started at the
median. 


We thus have a ``measurement'' of
position which disagrees with the Bohmian position.\footnote{Another 
example, in a certain sense simpler, is the two-dimensional
oscillator. For that example, the wave function\ evolution is still periodic in
time, but the different trajectories do not have a common period. See
\cite{DGZ03} for details.}


\subsection{Measurements of ``spin''}

Measurements of spin are the canonical measurements to discuss in
foundational matters. The basic experiment involves the Stern-Gerlach
device. Essentially, a  neutral spin-$\frac{1}{2}$ particle is shot into a
region with an inhomogeneous magnetic field. This magnetic field
interacts with the spin of the particle and forces the particle to
move up or down depending on whether the spin is up or down in the
appropriate direction. This is a repeatable
experiment. And, as is well known, a particle that has been measured spin up in
the $x$ direction and then spin up in the $z$ direction and is then
measured in the spin $x$ direction, will have spin up only 50\% of
the time. 

Bohmian mechanics adds just a little to the picture. The Stern-Gerlach
device splits the wave function\ into two pieces.  The particle travels in the
support of just one piece. If a measurement of the position of this
particle is made, then the conditional wave function\ changes; the ``empty''
packet is discarded. 
But if the system is kept
isolated, then the conditional wave function\ remains effective. In particular,
one might be able to arrange it so that the two pieces recombine. 
This is why
collapse can be a tricky concept and best discussed in the context of
Bohmian mechanics, a theory in which collapse 
can be analyzed. 

The most interesting aspect of the Bohmian story of spin measurements
is the following, as suggested in \cite{A92}. Prepare a system  with
sufficient symmetry and the particle will go up or down based on
whether it is above or below the line of symmetry. This is due to
the fact that Bohmian mechanics is a first-order theory.
One can setup the
Stern-Gerlach device so that going up registers as the particle being
spin up in the $z$ direction. One can also arrange the magnets, e.g.\
reverse the polarities, so that going up
represents being in the spin down state of the $z$ direction. Thus,
depending on the setup of the magnets, the very same Bohmian initial
conditions for the system
will give different results for the two experiments. This 
implies that the spin observable 
is not an intrinsic property of the system; rather, the spin
observable intimately depends on
the details of the experiment being done. That is essentially why the
theorems proving Bohmian mechanics impossible do not work. The
theorems assume  a
naive realist view towards the operators. But the operators are just
compact descriptions of the probability measure for the outcomes of
experiments. There does not have to be an underlying quantity common
to experiments measuring the same ``observable''.


\section{Global existence and uniqueness}

There are global existence results for complex-valued wave functions where the
configuration space is an open subset of \ensuremath{\mathbb{R}^{n}}.  
Although the
dynamics is undefined at the nodes, 
 it has been shown that typically the particle
evolution does not reach the nodes.  Indeed, it was shown 
in \cite{BDGZ95} that 
there is  
 global existence in time for typical 
 initial conditions where typicality is with respect 
to the measure $|\ensuremath{\psi}|^2$. This includes showing that the particles do
not reach, in finite time,  the nodes, the singularities of the
potential, or
 spatial
infinity.  

The proof is rooted in equivariance although equivariance only holds
once the existence results are established. Nevertheless, 
the basic idea is that the density vanishes at the places where the
particle dynamics ceases to exist. Since particles flow with the
density, this suggests that only a set of measure zero will reach the nodes,
singularities or infinity. The proof requires a series of estimates
involving the flux around these points. That typicality is all that
can be expected is demonstrated with a variety of examples.
These results do need to be extended to a more general setting, but
the expectation is that this is just a straightforward extension of
the arguments. 

\section{The historical story, hidden variables, and nonlocality}

Here is a rough sketch of the history. As I am not a qualified
historian, nor did I do any personal research into the story, please
take it only as my personal viewpoint. 

Bohmian mechanics has its roots at the very start of quantum
theory. Indeed, the folklore is that Einstein originally envisioned
the electromagnetic wave as guiding the particulate photons. Although
this came to nothing, it is not that distant from the Bohmian
viewpoint.  Bohmian mechanics was presented by Louis
de Broglie to the distinguished physicists of the Solovay
Congress. The lore has it that de Broglie's model was a bit overly
complicated and that he did not answer a question from Pauli in a
satisfactory manner. Thus, they dismissed the
theory. Also around this time, Madelung came up with a hydrodynamic
picture which chemists enjoy citing. It was also pretty close to Bohmian
mechanics. Nevertheless, the idea did not take hold. 

Instead, the quantum orthodoxy decided that such attempts were
impossible. They called such models hidden variable models. The idea
is that the state of the system is the wave function, and that  any
attempt to 
add additional variables, i.e.\ the hidden variables, in 
order to deduce the
measurement axioms  is doomed to failure. The first ``proofs'' were
given by von Neumann in the early 1930's. As usual, the mathematics
was fine, but the assumptions of what a hidden variable
theory had to satisfy were too restrictive. 

In the early 1950's, Bohm wrote a book on quantum theory, a book which
is still cited today as a good quantum book. In his book, he presented
a ``no hidden variables'' argument. Rumor has it that 
Bohm said that Einstein had convinced Bohm that
Bohm's argument did not work. His faith shaken in the standard dogma, 
Bohm came up with a counterexample to the ``no hidden variables''
theorem of von Neumann which we now call Bohmian mechanics. 
 Excitedly, he published his work. The physics community 
responded  negatively, but not convincingly. 


One mistake that Bohm made was this emphasis on the quantum
potential. This was an attempt to make the theory appeal to classical
intuitions. But as Bohmian mechanics is not at all a classical theory,
this made the theory rather complicated and it was deemed to be ad
hoc. Nevertheless, the physics community now had a counterexample to
the ``no hidden variables'' theorems. Instead 
of resolving the contradictions,  the physicists not only ignored
it, but continued to teach students that the type of theory that Bohm
had formulated was necessarily incompatible with quantum predictions.


Indeed, it seems that even Bohm abandoned his theory for a long while. It lay
dormant until J. S. Bell came across Bohm's theory. This was a
revelation to Bell, resulting in one of the most shocking 
theoretical and experimental results of the latter half of the \ensuremath{{20}^{th}}\
century. When Bell saw Bohm's theory, he realized two things. One was
the simple fact that it worked. Rather than focussing on the
              quantum potential,  he focussed on the density's being
equivariant. This allowed him to easily formulate the theory for
spin. 
But, more importantly, he realized that the theory was
nonlocal. This is obvious as the wave function is defined on the
configuration space, not on physical space. In fact, when evaluating 
the Bohmian velocity of a particle, one must, in general,
 take into account
 the positions of all the particles which seems inherently nonlocal. 
Bell 
asked whether one could remove the nonlocality. He also asked what
went wrong with the hidden variables arguments. 

The result of these considerations was Bell's famous inequalities. They
pertain to Bohm's version of the EPR experiment
 spins are used rather than position and momentum. 
The essential idea is that particles can be setup in such a way that
the measurement on one particle tells us the value of a relevant
observable for the other particle; the correlation is verifiable. 
The particles are to be widely
separated so that no local influences can intervene. By measuring
an incompatible observable on the second particle, one has results for
two incompatible observables. With only local influences propagating,
one would conclude that the values must already be predetermined. 
Bell's inequalities then showed that it is impossible for such values
to be predetermined assuming the quantum mechanical predictions. 
 The conclusion is not that hidden variables are
impossible. Rather, quantum mechanics is, fundamentally, nonlocal. To be sure,
the nonlocality seems benign enough to avoid signalling paradoxes, but
it is still true that the fundamental theory seems to be nonlocal. One
of the advantages of Bohmian mechanics is that it forces one to deal
with nonlocality in an honest way. 

Sadly, Bell's work was not a clarion call to embrace Bohmian
mechanics. Although Bell seems to have taken Bohmian mechanics
as the sensible version of quantum theory, he was invariably
misunderstood by the physics community. They would cite his work as
proof that hidden variables could not explain the results of
quantum measurements. His work actually shows that no local theory 
can explain it. But the point of particle motion is neither to
have locality nor to eliminate probabilities. Rather the point of
Bohmian mechanics is that it avoids the mysticism of quantum
mechanics.  For example, instead of postulating collapse,


one derives 
 collapse of the wave function\ from the theory. 


\section{Objections to Bohmian mechanics}

It seems that the objections to Bohmian mechanics come in four
generic varieties. The first is a misunderstanding of the facts of the
situation. To put it straight, Bohmian mechanics is a consistent
theory whose predictions agree with those of quantum mechanics. There
can be no disagreement. Thus, any evidence in support of orthodox
quantum theory is evidence for support of Bohmian mechanics.

This naturally leads to ``Who cares?'' There is not much to
say about this. We have already explained the clarity that Bohmian
mechanics gives to the measurement problem. That itself should be
sufficient. 

I would add that standard quantum mechanics appears as pure magic while 
Bohmian mechanics is the mechanism behind the show. The relevance of
this statement is in the teaching of quantum mechanics. The standard
approach mystifies students beyond the common mystification associated
with abstract mathematics. Fundamentally, many thoughtful 
students are not comfortable
with the measurement axioms. Unlike other areas of physics in which
explanations are offered,  quantum mechanics lends no explanations to
collapse or the special role of measurement. Adopting a Bohmian point
of view leads to a clear picture and, at the very least, one can cite
that the axioms mathematically follow from the evolution equations.  


There are many
objections to standard quantum theory that one can raise, but the
problem is that there is no actual standard quantum theory when
measurements are discussed. And it is exactly at the point of
measurements, at the contact with reality, that Bohmian mechanics is
relevant. It explains, in an obvious manner, what this mathematical
theory has to do with reality. More fundamentally, one should start,
when developing a theory, with an idea of what the primitive variables
are for the formulation of this theory. 
Primitive variables should not be a deduction of a theory,
although experiments may suggest suitable variables to choose. 
 Presumably the primitive variables should not be
results of measurements which are rather complicated objects centered
upon human existence. 
One does not want a human-centered theory; rather, one  wants
a clear theory in which certain behavior of the variables would give a
reasonable explanation of why the world behaves as it does. 

Bohmian mechanics postulates the existence of particles with
definite positions which evolve according to a prescribed dynamics. It happens
that this dynamics involves a wave function evolving according to
Schr\"odinger's equation. As we shall now explain, demanding a particle ontology almost
immediately leads to Bohmian mechanics.


\subsection{Where does Bohmian mechanics come from?}
The third objection to Bohmian mechanics is that it is cooked up. We
shall state five different derivations of Bohmian mechanics.


\begin{enumerate}
\item The starting place for deriving 
Bohm's equation is the same as for deriving
 Schr\"odinger's equation. We start with the idea that a wave is
involved. In particular, we assume we have a generic plane wave, i.e.\
\begin{equation} \label{bie:gp}
 \ensuremath{\psi}(\ensuremath{\boldsymbol{q}}, t) \ensuremath{:=} e^{i(\ensuremath{\boldsymbol{k}}
\cdot \ensuremath{\boldsymbol{q}} - \omega t)}.
\end{equation}
The empirical data somehow suggested that the momentum of the particle
was given by the de Broglie relation
\begin{equation} \label{bie:pk}
\ensuremath{\boldsymbol{p}} = \ensuremath{\hbar} \ensuremath{\boldsymbol{k}} . 
\end{equation}
The energy of the particle was related to the frequency via the Planck
formula 
\begin{equation} \label{bie:ew}
 E = \hbar \omega.
\end{equation}
We now look for an equation of motion for the particle as well as for the wave
equation. 

To obtain $\ensuremath{\boldsymbol{k}}$ from $\ensuremath{\psi}$, we take the spatial derivative
and divide by $i$ and $\ensuremath{\psi}$. Thus, \eqref{bie:pk} and \eqref{bie:gp}
suggests
that 
\[ m \ensuremath{\boldsymbol{v}} = \ensuremath{\boldsymbol{p}} = \ensuremath{\hbar} \ensuremath{\boldsymbol{k}} = \frac{\ensuremath{\hbar}}{i} \frac{\nabla \ensuremath{\psi}}{\ensuremath{\psi}}.\]
Considering, instead of \eqref{bie:gp}, a general complex-valued wave function, 
we find that we would end up
with a complex velocity vector. The simplest choice is to take the real
part of the expression above which, since we divided by $i$,
translates into taking the imaginary part. 
We have thus derived Bohm's equation:
\[
\ensuremath{\boldsymbol{v}}^{\ensuremath{\psi}} \ensuremath{:=} \frac{\ensuremath{\hbar}}{m} \textrm{Im} \frac{\nabla \ensuremath{\psi}}{\ensuremath{\psi}}.
\]

What about Schr\"odinger's equation? Well, classically, 
the energy of a particle is 
\[
E = \frac{\ensuremath{\boldsymbol{p}} \cdot \ensuremath{\boldsymbol{p}}}{2m} + V \label{bie:en} .
\]
We look for the simplest wave equation whose plane wave solutions
\eqref{bie:gp} are consistent with \eqref{bie:pk}, \eqref{bie:ew}, and
\eqref{bie:en}. 
We find that \eqref{bie:ew} implies for plane waves, that  
\[E = -\frac{\hbar}{i\ensuremath{\psi}} 
\ensuremath{\frac{\partial {\ensuremath{\psi}}}{\partial {t}}}.\] 
As for the $\ensuremath{\boldsymbol{p}}$ term, we have two possibilities:
\begin{align} 
\ensuremath{\boldsymbol{p}}
\cdot  \ensuremath{\boldsymbol{p}} =&  - \hbar^2
\frac{\nabla \ensuremath{\psi}}{\ensuremath{\psi}} \cdot \frac{\nabla \ensuremath{\psi}}{\ensuremath{\psi}}  \label{bie:p2a}
\end{align}
or 
\begin{align}
\label{bie:p2b}
 \ensuremath{\boldsymbol{p}}
\cdot  \ensuremath{\boldsymbol{p}} =&
-\hbar^2 \frac{\ensuremath{\Delta} \ensuremath{\psi}}{\ensuremath{\psi}}.
\end{align}
Choosing \eqref{bie:p2b}
 leads to  
\[
-\frac{\hbar}{i\ensuremath{\psi}} 
\ensuremath{\frac{\partial {\ensuremath{\psi}}}{\partial {t}}} = -\frac{\hbar^2}{2m} \frac{\ensuremath{\Delta} \ensuremath{\psi}}{\ensuremath{\psi}} +V,
\]
which becomes Schr\"odinger's wave equation after multiplying through by $\ensuremath{\psi}$.
As for \eqref{bie:p2a}, this is not a linear equation. For a linear
equation, knowing the solutions for any initial plane wave actually
 specifies the
evolution for any function. This is not true for nonlinear equations. 

We have arrived at Bohmian mechanics. We shall discuss in the companion manuscript on geometrically formulating Bohmian theories how to arrive at the Pauli equation. 

\item Briefly, if we start with Schr\"odinger's equation, then we can derive a
continuity equation for $\ensuremath{\rho} = \ensuremath{|\ensuremath{\psi}|^2}$. By dividing the probability
current by $\ensuremath{\rho}$, we obtain a velocity field, 
the Bohm velocity field. Equivariance is then immediate, and with it
the 
empirical agreement with quantum mechanics. In terms of
finding Bohmian laws of motion, this has been a successful and easy
method for determining appropriate laws.
 But as a derivation, it has a
feeling of fitting theory to given data. Thus, it is best if we have
other methods. But we do note that 
 one of the guiding principles in forming quantum
theories is to establish that there is an appropriate 
probability current associated with the equation. This suggests that
we should perhaps formulate the Bohmian law of motion and then find a
wave equation based on demanding equivariance.

\item Schr\"odinger's equation also implies a Hamilton-Jacobi equation. This
is Bohm's method of derivation. It
describes particles moving under a force consisting of the classical
force and a new quantum force. 
But we have to
supplement the force with the constraint on the initial velocity field 
 given by
the wave function according to Bohm's equation. The strange feature
about this is that the constraint is actually the law of motion one
wants. 


\item In \cite{DGZ92}, the authors give a rather convincing derivation
of Bohm's equation based on symmetry considerations. The quantity 
$\nabla \ensuremath{\psi}$ is a natural quantity for a rotation invariant way of
obtaining a vector. Dividing by \ensuremath{\psi}\ comes from noting that it is only
the ray in Hilbert space which is physically relevant. 
The imaginary part arises from time reversal invariance. Under time
reversal, the wave function\ is sent to its complex conjugate and velocity
fields change sign. 
To obtain the
factor of $\frac{\ensuremath{\hbar}}{m}$, we must consider Galilean boosts. Under a
boost, we have that $\ensuremath{\psi} \mapsto e^{-i\frac{m}{\ensuremath{\hbar}} \ensuremath{\boldsymbol{w}} \cdot \ensuremath{\boldsymbol{q}}} \ensuremath{\psi}$ 
where $\ensuremath{\boldsymbol{w}}$ is the velocity of the new frame with respect to the old frame. A
velocity $\ensuremath{\boldsymbol{v}}$ should tansform to $\ensuremath{\boldsymbol{v}} - \ensuremath{\boldsymbol{w}}$ in the new
frame and the
Bohm velocity will become $\ensuremath{\boldsymbol{v}}^{\ensuremath{\psi}}_{old} - \ensuremath{\boldsymbol{w}}$ if and only if the
factor in front is $\frac{\ensuremath{\hbar}}{m}$. 

\item A new derivation is in \cite{DGTuZ03a}. It seems to apply to a
  wide variety of cases, such as the non-relativistic case, 
quantum field theory, and the Dirac equation. In some situations, it
must be interpreted correctly. We shall only show here that it agrees
in the non-relativistic case. The idea is that a tangent vector is a
first-order differential operator acting on functions and obeying the
Leibniz rule. Thus, $v$ is defined if we know what $v(f)$ is for all
$f$. In particular, if we have coordinates, then the tangent vector
$v$ has coordinate components $v^j = v(x^j)$. We also note that if
$q(t)$ is a trajectory whose tangent vector is $v$ at time $0$, then
defining $g \ensuremath{:=} f \circ q$, we have
$\dot{g} (0) = v(f)$. 

From a quantum perspective, functions can be thought of as
multiplication operators acting on the Hilbert space of wave functions. In
general, an operator $A$'s time evolution, in the Heisenberg picture, is
given by\footnote{
This follows by shifting from the Schr\"odinger\ picture, in which $A$ does
not evolve and the expected value of the operator at time $t$ is given
by \ensuremath{({\ensuremath{\psi}_t}, {A \ensuremath{\psi}_t})}, to the Heisenberg picture in which only the
operator evolves and the expected value at time $t$ is \ensuremath{({\ensuremath{\psi}_0}, {A_t
\ensuremath{\psi}_0})}. Indeed, let $U_t = e^{-\frac{i}{\hbar}Ht}$ be the (unitary) Schr\"odinger\
evolution operator. Then to 
obtain the same answer in both cases, we need $A_t = U_t^{-1} A
U_t$.
By differentiating with respect to time, we find that 
\[\dot{A}_t =
\frac{i}{\hbar} H A_t + A_t (-\frac{i}{\hbar} H)= \frac{i}{\hbar}
[H,A_t].\]
} 
\[ \dot{A}_t = \frac{i}{\ensuremath{\hbar}} [H, A_t].
 \]
Thus, transform $f$ into a multiplication operator: $(\hat{f} \ensuremath{\psi}
)(q) \ensuremath{:=} f(q) \ensuremath{\psi}(q).$ We shall assume that $f$ is smooth and is
compactly supported, as local information is all that we care about. Then
the Bohmian velocity field $v$ with a given $\ensuremath{\psi}$ 
can be defined as the differential
operator whose action satisfies
\[ 
(v (f))(q) = \dot{f} (q) = \textrm{Re}
\frac{\ensuremath{({\ensuremath{\psi}}, {\dot{\hat{f}} \ensuremath{\psi}})}}{\ensuremath{({\ensuremath{\psi}}, {\ensuremath{\psi}})}} (q) 
=  \textrm{Re} 
\frac{\ensuremath{({\ensuremath{\psi}}, {\frac{i}{\hbar} [H,\hat f] \ensuremath{\psi}})}}{\ensuremath{({\ensuremath{\psi}}, {\ensuremath{\psi}})}} (q).
\]

We shall now investigate this for Schr\"odinger's equation. First, note that any
potentials in $H$ commute with $f$ and thus they contribute nothing to
the velocity field. Let us therefore take $H$ to be $-\frac{\ensuremath{\hbar}^2}{2m}
\ensuremath{\Delta}$. 
The commutator is easy to compute; we find
\[
[\ensuremath{\Delta}, \hat{f}] \ensuremath{\psi} =  \ensuremath{\Delta} f \ensuremath{\psi} + 2 \nabla \ensuremath{\psi} \cdot \nabla f . 
\]
Thus, we have that 
\[
v(f) = v \cdot \nabla f
\]
where $v$ is the Bohm velocity field.


\end{enumerate}

\subsection{Relativistic extensions of Bohmian mechanics}
The fourth objection is that it is irreconcilable with quantum field
theory and with relativity. These claims are false as recent work has
demonstrated. 

Nonlocality is often assumed to be problematical for a hidden variable
perspective. There are two responses to this. As stated previously, it is a
problem for standard quantum theory as well, but due to the shiftiness
of what is being described, the problem can be hidden for a long
time. The second point is that there is no inherent inconsistency with
nonlocal influences in a relativistic theory, as discussed at length
in \cite{M94}. More to the point, there are relativistic,
nonlocal Bohmian models \cite{BDGZ99, GTu03}.\footnote{One might
  argue as to what Bohmian means; in this section, we take it to mean
  a theory in which there is a clear ontology with prescribed dynamics
  and no reference to measurements in the formulation of the theory.}
The configuration space for $N$ particles moving in the space-time
manifold $M$ might
be taken to be $M^N$. Then  
the wave function\ in a relativistic setting
would seem to be a function on $M^N$. That
is to say, it is a  multi-time wave function. When
evaluating the Bohmian evolution, one must evaluate the wave function\ by using
the 
positions of the other particles. The problem is that we no longer
have simultaneity to tell us what the full configuration is. 
One model involves a foliation of space-time. That is to say,
simultaneity surfaces are built into the structure. Although perhaps
distasteful, if a compelling Lorentz-invariant 
law for the foliation can be found, then this should be
acceptable. With the foliation method, the empirical predictions are
transparent and they are equivalent to those of quantum theory. In
particular, the foliation is unobservable. Another idea involves no
additional structure. It uses the light cones to accomplish what the
simultaneity surfaces do. To have a chance at agreeing with
experiments, it is necessary to use the forward light cones. Thus, one
might say there are two different arrows of time involved. But 
the probabilistic import of the model is still unclear. Existence of
the dynamics is also rather unclear.


Quantum field theory is perhaps the object most pointed towards when
decrying the Bohmian point of view. Bohm took the approach of giving
ontological status to the fields. A recent approach \cite{DGTuZ02} is to retain
particles as the ontology with the wave function\ retaining its role as guiding
the particles. Describing the evolution of the wave function\ is the primary
concern of quantum field theory and it is difficult.  Understanding
the role of particles is quite easy. The wave function\ is a function on the
disjoint union of all of the configuration spaces for the various
number of particles. That is, there is no fixed number of
particles. The Bohmian motion consists of two parts. One is the usual
part in which the particles evolve deterministically guided 
the wave function. 
The other part is stochastic. In a random fashion, particles
are created or annihilated. The motion
proposed is equivariant implying that this model
gives the same empirical 
predictions as standard quantum field theory. But it is
a theory in which no discussion about measurements is necessary in
order to make this a theory about reality. It is true that this is not
deterministic, but that was never the issue driving Bohmian
mechanics. The real issue was the demand for a simple,
well-defined theory of reality.

\clearpage
\bibliographystyle{abbrv}
\bibliography{references}
\end{document}
